\documentclass[11pt]{article}
\usepackage[margin=1in]{geometry}
\usepackage{times}
\usepackage{amsmath}
\usepackage{amssymb}
\usepackage{graphicx}
\usepackage{booktabs}
\usepackage{array}
\usepackage{hyperref}
\usepackage{natbib}
\usepackage{enumitem}
\usepackage{caption}
\hypersetup{colorlinks=true, linkcolor=blue, citecolor=blue, urlcolor=blue}

\title{\textbf{Behavioral Control Through System Prompts:\\
A Comparative Structural Analysis of Public Governance Artifacts\\
Across Commercial LLM Deployments}}

\author{Syed Abdul Aman \\
Independent Researcher \\
\texttt{amanbaba9404522@gmail.com}}

\date{July 2026}

\begin{document}

\maketitle

\begin{abstract}
Large language model providers increasingly publish documents intended to govern
model behavior at scale, ranging from training-time value specifications to
runtime authority hierarchies to per-product deployed instructions. Despite
substantial public interest in these artifacts --- evidenced by prompt-collection
repositories with tens of thousands of stars and an active empirical literature
on prompt-extraction attacks \citep{promptleak2026} --- direct comparative content
analysis of these documents' primary text remains rare. Prior
literature surveys task-oriented prompting techniques for benchmarking
(which concerns user-authored prompts, not provider-authored governance),
studies the security of prompt confidentiality (which concerns extraction,
not content), or examines system-level instructions as an object of research
and policy discourse \citep{neumann2026governance} without directly coding
the primary governance documents themselves. This paper
introduces a three-part taxonomy of behavioral-control artifacts --- training-time
constitutions, runtime authority specifications, and deployed inference-time
system prompts --- and applies a credibility-tiered comparative methodology
across four major labs (Anthropic, OpenAI, Google DeepMind, xAI) and a
25-plus-product corpus of coding-agent tools. We find that labs diverge not
only in the content of behavioral constraints but in which artifact type they
use to express them at all; that transparency is sometimes reactive
(incident-driven) rather than proactive; that product category (consumer chat
versus coding agent) systematically predicts which instruction categories
dominate a deployed prompt; and, using a publicly verifiable 18-revision
changelog spanning 23 months, that at least one lab's self-description of its
governance artifact as a ``living document'' can be independently confirmed
against dated evidence rather than taken on trust. We argue this taxonomy
resolves a category error common in informal prompt-comparison work, and we
outline a research agenda for treating behavioral-control artifacts as
first-class objects of study in AI governance and NLP.
\end{abstract}

\noindent\textbf{Keywords:} system prompts, LLM alignment, AI governance,
prompt engineering, behavioral control, transparency

\section{Introduction}

Every deployed conversational AI system operates under some form of
provider-authored instruction layer that sits above the user's own input ---
commonly called a ``system prompt.'' As these systems have moved from research
demonstrations to consumer and enterprise products, the documents that
constrain their behavior have grown from short, informal instructions into
substantial artifacts: Anthropic's constitution for Claude runs to roughly
23,000 words \citep{anthropic2026constitution}; OpenAI's Model Spec runs to
128 pages in the revision analyzed here \citep{openai2025modelspec}.\footnote{We
analyze the 2025-09-12 revision throughout this paper. OpenAI has continued
to revise the Model Spec since (including further updates through at least
March 2026); as with all artifacts in this study, we report a dated snapshot
rather than a currently-synchronized document, consistent with the
credibility-tiering and snapshot-transparency approach described in
Section~\ref{sec:methodology}.} At the same time, a large ecosystem of
public repositories exists purely to collect and share these documents,
whether officially released or reconstructed through extraction, with the
single largest such repository exceeding 140,000 stars on GitHub at the time
of writing \citep{x1xhlol2026repo}.

Despite this level of public attention, the documents themselves have
rarely been treated as a coherent object of direct comparative content
analysis. Three adjacent literatures exist. The first is the
prompt-engineering literature, which develops taxonomies of \emph{task}
prompts --- the instructions a user or application developer supplies to
elicit a specific output --- for purposes of benchmarking and technique
comparison \citep{teler2023,liu2024survey}. The second is the AI security
literature, which studies system prompts as a \emph{target}: how they can
be extracted via adversarial querying, and how such leakage can be
mitigated \citep{promptleak2026,proxyprompt2025}. The third, most closely
related to this paper, is a contemporaneous governance-studies literature
exemplified by \citet{neumann2026governance}, which systematically reviews
how researchers and policymakers \emph{discuss} system-level instructions
as governance objects, but does not itself perform a direct, coded
comparison of the documents' primary text across labs. None of these three
literatures asks the question this paper asks: across labs and product
categories, what do these documents actually contain, how are they
structured, and what does the variation reveal about differing approaches
to AI governance?

This gap matters for at least three reasons. First, from a governance
perspective, these documents are among the only concrete, inspectable
artifacts through which the public can observe how a lab operationalizes
its stated values --- more concrete than a blog post, more current than a
research paper. Second, from an NLP research perspective, treating
governance documents as training or evaluation data (as some labs already
do internally) requires a shared vocabulary for what these documents
contain, which does not yet exist in the public literature. Third, from a
practical engineering perspective, developers building on top of these
models routinely conflate categorically different artifacts --- for
instance, comparing a leaked \emph{deployed} prompt for one product against
an official \emph{training-time} document for another --- producing
comparisons that are not actually commensurable.

\subsection*{Contributions}
\begin{enumerate}[leftmargin=*]
  \item We introduce a three-part taxonomy distinguishing training-time
    constitutions, runtime authority specifications, and deployed
    inference-time system prompts as categorically distinct artifact types,
    and show that failing to distinguish them produces invalid comparisons.
  \item We construct a credibility-tiered comparative dataset spanning
    four major labs and a broad corpus of coding-agent products, coded
    along eight functional axes (identity framing, safety/refusal
    instructions, tool-use governance, tone/persona control, output format
    constraints, injection countermeasures, recency handling, and length).
  \item We report empirical findings on cross-lab divergence, including an
    asymmetry in whether labs publish governance artifacts at all, a case
    where transparency was reactive rather than proactive, a systematic
    difference in instruction emphasis between consumer-chat and
    coding-agent product categories, and --- using a publicly available,
    dated changelog --- a directly verifiable measurement of publication
    cadence for one lab's deployed system prompt (18 revisions over 23
    months), rather than relying on a publisher's self-description of its
    artifact as a ``living document.''
  \item We propose a research agenda for treating behavioral-control
    artifacts as first-class objects of study.
\end{enumerate}

\section{Related Work}

\subsection{Task-Prompt Taxonomies}
The dominant strand of prompt-taxonomy research addresses how to structure
task instructions to maximize model performance on downstream benchmarks.
The TELeR taxonomy \citep{teler2023} categorizes prompts by turn structure,
expression style, level of detail, and role, explicitly for the purpose of
standardizing benchmark reporting. Broader surveys of prompt-engineering
technique \citep{liu2024survey} catalog methods such as chain-of-thought,
few-shot exemplification, and retrieval augmentation. A related body of work
constructs taxonomies of \emph{prompt defects} --- failure modes in
task-prompt design drawn from software-engineering and NLP venues
\citep{promptdefects2025}. All of this work treats the prompt as something
a user or developer constructs to accomplish a task; none of it addresses
the provider-authored instruction layer that this paper studies.

\subsection{Prompt Security and Extraction}
A substantial and active literature studies system prompts as
confidentiality targets. Early work formalized prompt injection and prompt
leaking as distinct attack classes \citep{perez2022}. Subsequent work has
measured leakage empirically across large numbers of deployed applications,
finding leakage rates exceeding 80\% in some large-scale studies
\citep{promptleak2026}, and has proposed defenses such as proxy substitution
\citep{proxyprompt2025} and encoding-based hardening
\citep{encodingdefense2026}. This literature's object of study is the
\emph{event} of extraction and the \emph{robustness} of the prompt against
it, not the semantic content of the prompt as a governance artifact. Our
work is complementary: where security research asks whether a prompt can be
extracted, we ask what it says once available through legitimate or
well-established channels, and what patterns exist across many such
documents.

\subsection{AI Governance Transparency}
A smaller literature and body of industry practice addresses transparency
in AI governance more broadly, including model cards, system cards, and
published behavioral specifications. OpenAI's Model Spec explicitly
positions itself as a transparency artifact governing the ``chain of
command'' between the provider, operators (system-prompt authors), and
users \citep{openai2025modelspec}. Anthropic's constitution is explicitly
framed as a document written primarily \emph{to} the model itself, optimized
for precision over public accessibility, with public release framed as a
secondary transparency benefit \citep{anthropic2026constitution}.

The closest prior work to this paper is a contemporaneous study by
\citet{neumann2026governance}, published at FAccT 2026, which treats
system-level instructions as a governance object through a PRISMA-based
systematic literature review of 287 papers (yielding an eight-category
typology of the \emph{goals} researchers attribute to such instructions,
e.g., alignment, stability, auditability) combined with a policy-document
analysis of two case studies, the US Executive Order 14319 and the EU
General-Purpose AI Code of Practice. Their central finding --- that
governance frameworks often assume system-level instructions provide
stable, predictable behavioral control in ways not consistently supported
by the technical literature --- is a meta-level claim about the literature
and policy discourse \emph{surrounding} these artifacts. Their organizing
axis is the \emph{stakeholder hierarchy} through which instructions are
issued (a ``prompt stack'' of foundation-model provider, developer/deployer,
and end user). Our study differs in scope and method rather than
contradicting theirs: we do not survey what researchers or policymakers
claim about system-level instructions, but instead directly and
comparatively code the content of the primary governance documents
themselves, organized around a taxonomy of \emph{artifact type}
(training-time, runtime-authority, and deployed) that is orthogonal to
the stakeholder-hierarchy axis used in their typology. Notably, based on
the sections and reference list available to us at the time of writing,
Neumann et al.'s corpus draws on Anthropic's published deployed system
prompts but does not appear to treat Anthropic's training-time constitution
as a separately governed artifact type, which is consistent with our
taxonomy identifying it as a categorically distinct object from the
runtime and deployed layers their stakeholder-hierarchy framing is best
suited to describe. We view the two studies as complementary: theirs
establishes, at the level of the research and policy literature, that the
premise of prompt-based governance is more contested than commonly
assumed; ours provides a primary-source comparative account of what the
documents underlying that premise actually contain.

\section{A Taxonomy of Behavioral-Control Artifacts}
\label{sec:taxonomy}

Informal comparisons of ``system prompts'' across labs routinely conflate
three categorically distinct artifacts. We propose the following taxonomy,
motivated directly by the sources examined in this study.

\begin{description}[itemsep=6pt]
  \item[Training-time constitutions.] Documents used during model training
    to shape weights, typically via constitutional-AI-style methods in which
    the document generates synthetic training signal (example interactions,
    response rankings, scenario guidance). Anthropic's constitution is the
    clearest public instance: it is addressed to Claude as its primary
    audience, is not injected at inference time, and explicitly avoids
    prescribing rule-like behavior in favor of reasoning the model is
    expected to internalize \citep{anthropic2026constitution}.
  \item[Runtime authority specifications.] Documents that govern how a
    deployed model should arbitrate between instructions from different
    sources at inference time, without themselves being a single prompt.
    OpenAI's Model Spec is the clearest instance: it defines a ranked
    ``principal hierarchy'' (provider, operator, user) and explains how
    conflicts between levels should resolve, but is not itself the text
    injected into any single conversation \citep{openai2025modelspec}.
  \item[Deployed inference-time system prompts.] The specific text actually
    prepended to a conversation for a given product at a given time. This
    is the artifact most commonly leaked, extracted, or --- in xAI's case
    \citep{xaigrokrepo2026} --- proactively published in raw form. It is the
    most ephemeral of the three, changing frequently with product updates,
    and the most heterogeneous, since a single lab may maintain dozens of
    distinct deployed prompts across product surfaces.
\end{description}

This taxonomy is not merely descriptive bookkeeping. Section~\ref{sec:findings}
shows that failing to distinguish these types produces comparisons that
are not commensurable, and that the choice of \emph{which} artifact type a
lab uses to express a given constraint is itself informative about that
lab's governance philosophy.

\section{Methodology}
\label{sec:methodology}

\subsection{Source Tiering}
Because the corpus necessarily includes documents of varying provenance, each
entry is tagged with a credibility tier, reported transparently rather than
treated as uniformly authoritative:

\begin{itemize}[leftmargin=*]
  \item \textbf{Tier 1 (Official/primary):} Content published directly by the
    originating lab (e.g., Anthropic's constitution, OpenAI's Model Spec,
    xAI's raw prompt repository).
  \item \textbf{Tier 2 (Reconstructed, high-consensus):} Content appearing
    consistently across independent extraction attempts, hosted in
    long-running, widely-cited public repositories (e.g., the
    coding-agent prompt aggregation studied in
    Section~\ref{sec:coding-agents}).
  \item \textbf{Tier 3 (Single-source leaks):} Lower-confidence material used
    only to note the possible existence of a pattern, never as a primary
    data point.
\end{itemize}

\subsection{Coding Scheme}
Each entry in the corpus is coded along eight functional axes: identity
framing, safety/refusal instructions, tool-use governance, tone/persona
control, output format constraints, injection/jailbreak countermeasures,
recency/knowledge-cutoff handling, and document length. This scheme was
derived inductively from an initial pass over the corpus rather than
imposed a priori, consistent with standard qualitative coding practice for
emergent taxonomies \citep{promptdefects2025}.

\subsection{Copyright and Reproduction Policy}
No system prompt text is reproduced verbatim beyond short quoted fragments
where exact wording is analytically load-bearing. All classification in
this paper is performed through structural coding and paraphrase, not
reproduction, consistent with treating these as evidentiary rather than
literary objects.

\subsection{Scope and Corpus}
The comparative corpus (Table~\ref{tab:corpus}) spans four major labs at
the level of official or near-official governance artifacts, and a
25-plus-product corpus of coding-agent system prompts aggregated from a
single long-running, high-consensus Tier-2 source
\citep{x1xhlol2026repo}, used to characterize product-category-level
patterns rather than individual product claims.

\begin{table}[htbp]
\centering
\small
\begin{tabular}{@{}p{2.2cm}p{2.9cm}p{1cm}p{6.2cm}@{}}
\toprule
\textbf{Lab} & \textbf{Artifact} & \textbf{Tier} & \textbf{Type (per Sec.~\ref{sec:taxonomy})} \\
\midrule
Anthropic & Constitution & 1 & Training-time \\
Anthropic & \texttt{claude.ai}/mobile system-prompt changelog & 1 & Deployed (versioned, 18 dated revisions, Jul.\ 2024--Jun.\ 2026) \\
OpenAI & Model Spec & 1 & Runtime authority spec \\
xAI & \texttt{grok-prompts} repo & 1 & Deployed (multiple, raw) \\
Google DeepMind & Prompting guides & 3 & None published (developer-facing only) \\
25+ coding-agent products & Aggregated repo & 2 & Deployed (per-product) \\
\bottomrule
\end{tabular}
\caption{Corpus overview by lab/product category, credibility tier, and
artifact type. Anthropic is the only lab in the corpus with both a
training-time and a versioned, longitudinally documented deployed artifact
publicly available.}
\label{tab:corpus}
\end{table}

\section{Findings}
\label{sec:findings}

\subsection{Labs Diverge in Artifact \emph{Type}, Not Just Content}
The most basic finding is that labs do not agree on what kind of document
should perform behavioral governance. Anthropic's primary public artifact is
a training-time document not injected at runtime. OpenAI's primary public
artifact is a runtime arbitration framework, not itself a prompt. xAI's
primary public artifact is the deployed prompts themselves, in raw form,
with no equivalent training-time or arbitration-framework document made
public. Google DeepMind has published no artifact in any of the three
categories that addresses the company's own behavioral constraints on
Gemini; its public documentation is prompting guidance addressed to
third-party developers, not a transparency document about Gemini's own
governance. A naive comparison of ``Anthropic's system prompt vs.\ OpenAI's
system prompt vs.\ Google's system prompt'' is therefore not comparing four
instances of the same artifact type, but at best two comparable instances
and two absences.

\subsection{Convergent Hierarchical Framing, Divergent Justification}
Despite this divergence in artifact type, Anthropic and OpenAI converge on
a strikingly similar structural device: an explicit, ranked hierarchy of
authority. OpenAI's principal hierarchy ranks provider, operator, and user
instructions, with higher levels overriding lower ones in conflict
\citep{openai2025modelspec}. Anthropic's constitution specifies a priority
order of safety, ethics, compliance, and helpfulness. The two labs differ
sharply, however, in how they justify the hierarchy: OpenAI frames it in
terms of steerability and developer control, drawing an analogy to
function definitions and arguments in a programming language, while
Anthropic explicitly frames its equivalent structure in terms of reasoning
the model should internalize rather than rules to mechanically apply,
and explicitly states a preference for reason-based over rule-based
specification on the grounds that rule lists do not generalize to novel
situations. This is a substantive difference in governance philosophy
expressed through a superficially similar structural device, and would be
invisible to an analysis that only asked ``does the document specify a
hierarchy?'' without examining the justificatory framing.

\subsection{Transparency Can Be Reactive, Not Only Proactive}
xAI's decision to publish Grok's raw system prompts is frequently described
in press coverage as a transparency initiative comparable to Anthropic's.
The provenance differs materially: xAI's public commitment to publish
prompts was announced as a direct response to a May 2025 incident in which
an unauthorized internal modification caused the Grok bot on the X platform
to produce off-topic, ideologically directed outputs \citep{xaiincident2025}.
Publication was explicitly framed as a trust-repair mechanism and paired
with a commitment to additional internal review controls, rather than
introduced as a proactive philosophical stance of the kind that
characterizes Anthropic's or OpenAI's framing. This suggests that
``publishes its system prompts'' is not a single governance stance but a
category that includes both prospective transparency-as-values and
reactive transparency-as-incident-response, with potentially different
implications for how much a given publication should be trusted as
representative of typical behavior going forward. Consistent with this, at
least one public issue thread on xAI's repository documents a case where a
prompt reportedly shared directly by the model in conversation was not
subsequently found in the public repository, raising unresolved questions
about repository completeness that this paper does not attempt to
adjudicate but flags as a limitation of Tier-1 status not implying
completeness.

\subsection{Substantive, Not Just Structural, Divergence}
Labs also diverge in the actual content of behavioral constraints in ways
that reflect differing product positioning rather than differing
communication style. xAI's published safety-prefix file for its
consumer-facing model states that no additional content restrictions apply
to fictional adult sexual content involving dark or violent themes unless
specified elsewhere, and separately instructs the model to treat users as
adults and avoid moralizing about ``edgy'' requests, with these
instructions marked as taking the highest priority over other
instructions \citep{xaigrokrepo2026}. This is a substantively more
permissive default stance on mature fictional content than is publicly
documented for the other labs in this corpus, while the same published
material maintains a narrower, explicitly named restriction on content
involving minors. The combination --- broad default permissiveness paired
with one hard, explicitly carved-out exception --- appears to be a
deliberate product-differentiation choice rather than an oversight, and
is stated plainly enough in the primary source that it requires no
inference on our part.

\subsection{Product Category Predicts Instruction Emphasis}
\label{sec:coding-agents}
Coding-agent products (Cursor, Claude Code, GitHub Copilot, Devin, and
over twenty others aggregated in a single high-consensus Tier-2 source)
show a systematic pattern distinct from consumer chat assistants: heavy
emphasis on tool-orchestration governance (sequencing of tool calls,
confirmation requirements before destructive file operations, diff-format
output conventions) and comparatively light emphasis on general
content-safety framing. We interpret this as consistent with the
three-artifact taxonomy in Section~\ref{sec:taxonomy}: where the underlying
base model already carries general safety behavior from a training-time
layer (a constitution or equivalent), the deployed, product-specific
system prompt's marginal job is to govern what is specific to that
product's task domain, not to re-specify general safety behavior the base
model already encodes. This is consistent with, though not proof of, a
layered division of governance labor across the three artifact types.

\subsection{Publication Cadence Reveals Living-Document Status Empirically, Not Just Rhetorically}
\label{sec:cadence}
Both the Anthropic constitution and the OpenAI Model Spec describe themselves in their own text as living documents subject to revision. Anthropic is the only lab in this corpus for which that claim can be verified against a public, dated record rather than taken on the publisher's word: its deployed \texttt{claude.ai}/mobile system-prompt changelog lists 18 distinct dated revision events across roughly 23 months (July 2024 to June 2026), a mean interval of approximately 41 days, or under six weeks, between documented changes \citep{anthropic2026sysprompts}. The changelog further states explicitly that these updates do not propagate to the API, confirming --- independent of any inference on our part --- that the deployed prompt is a product-surface-specific artifact rather than a single canonical text shared across all of a lab's surfaces, consistent with the taxonomy in Section~\ref{sec:taxonomy}. No equivalent longitudinal record is publicly available for OpenAI, Google DeepMind, or xAI at the time of writing, which limits comparison of update cadence across labs to this single case; we report it not as a cross-lab comparison but as an existence proof that at least one lab's ``living document'' framing is independently verifiable rather than solely rhetorical, and as a template for what verifiable transparency could look like elsewhere in the corpus.

\section{Discussion}

The findings above suggest two things. First, treating ``system prompt
comparison'' as a well-posed question requires first answering \emph{which
artifact type} is being compared, and being explicit when a comparison
spans types (e.g., a leaked deployed prompt against an official
training-time document) rather than presenting it as apples-to-apples.
Second, the variation observed across labs is not merely stylistic. The
choice of artifact type, the decision whether to publish anything at all,
and the circumstances under which publication occurs (proactive versus
incident-driven) are themselves meaningful data about a lab's governance
posture and should be treated as such by researchers and by the broader
public discourse that currently tends to flatten ``does lab X publish its
prompts'' into a single binary marker of trustworthiness.

\subsection{Implications for Practitioners}
Developers building on top of multiple providers' APIs, including within
agentic systems that call more than one model, should be aware that the
``safety behavior'' of a product is not uniformly encoded at the same
layer across providers, and that a system prompt inherited from one
product's convention (e.g., a coding-agent's typically light safety
framing) should not be assumed portable to a different deployment context
without re-examining what layer is expected to carry that responsibility.

\subsection{Implications for Researchers}
We suggest that behavioral-control artifacts be treated as a legitimate,
citable object of NLP and AI-governance research in their own right, with
the credibility-tiering approach used here as one candidate methodological
norm for handling the mix of official and reconstructed material that
necessarily characterizes this domain until more labs adopt full proactive
publication.

\section{Limitations}

This study has several limitations that should inform how its findings are
weighted. First, the corpus is not exhaustive: it covers four major labs
in depth and one aggregated coding-agent source, not the full space of
deployed LLM products. Second, Tier-2 and Tier-3 material, by construction,
carries lower evidentiary weight than Tier-1 material, and conclusions
drawn substantially from Tier-2 sources (notably the coding-agent
product-category finding in Section~\ref{sec:coding-agents}) should be
read as suggestive pending independent replication against official
sources where labs choose to publish them. Third, most artifacts studied
here are single or few-point snapshots rather than longitudinal series;
the constitution and Model Spec documents are both explicitly described
by their publishers as living documents subject to revision, but only
Anthropic's deployed-prompt changelog (Section~\ref{sec:cadence}) provides public,
dated, longitudinal evidence of that revision in practice. Findings about
relative emphasis or structure for OpenAI, Google DeepMind, and xAI may
therefore not hold at a later snapshot, and we have no public basis for
estimating how quickly they would change if they did. Fourth, the coding
scheme, while derived inductively from the corpus, was applied by a single
coder; inter-rater reliability was not assessed and would strengthen
future iterations of this work. Fifth, the interpretive claim in
Section~\ref{sec:coding-agents} that base-model training absorbs general
safety behavior while deployed prompts absorb task-specific governance is
consistent with the observed pattern but was not tested against a
condition that could falsify it (e.g., a coding agent built on a base
model with no equivalent training-time safety layer); we present it as
the most parsimonious explanation available from public evidence, not as
a demonstrated causal mechanism.

\section{Conclusion and Future Work}

Behavioral-control documents governing large language models have grown
from informal instructions into substantial public artifacts, yet direct
comparative content analysis of their primary text --- as distinct from
literature and policy review of how such documents are discussed
\citep{neumann2026governance} --- has remained rare. We
introduced a three-part taxonomy distinguishing training-time
constitutions, runtime authority specifications, and deployed
inference-time system prompts, and showed that this distinction resolves
otherwise incommensurable comparisons across labs. Applying a
credibility-tiered methodology across four major labs and a broad
coding-agent product corpus, we found that labs diverge in which artifact
type they use at all, that transparency can be reactive rather than
proactive, that content divergence reflects genuine product-positioning
choices rather than only communication style, that product category
systematically predicts which instruction categories a deployed prompt
emphasizes, and that at least one lab's claim to maintain a living
governance document is independently verifiable against a public,
dated record rather than resting on the publisher's word alone.

Future work should extend the corpus to additional labs and product
categories as more official material becomes available, assess inter-rater
reliability on the coding scheme, and track these artifacts longitudinally
given their explicitly living-document status, which would allow the
field to study not just cross-sectional divergence but the trajectory of
behavioral-control design over time.

\bibliographystyle{plainnat}
\begin{thebibliography}{99}

\bibitem[Anthropic(2026a)]{anthropic2026constitution}
Anthropic. \emph{Claude's Constitution}. January 2026.
\url{https://www.anthropic.com/constitution}

\bibitem[Anthropic(2026b)]{anthropic2026sysprompts}
Anthropic. \emph{System Prompts} [release-notes documentation, versioned
changelog of the deployed claude.ai/mobile system prompt, dated entries
from July 2024 through June 2026]. Accessed July 2026.
\url{https://docs.claude.com/en/release-notes/system-prompts}

\bibitem[Neumann et al.(2026)]{neumann2026governance}
A.~Neumann, H.~Sargeant, and J.~Singh. Prompt Governance? On Governing
Technologies Governed by Natural Language. \emph{arXiv:2606.07539}, 2026.
Also published in \emph{Proceedings of the 2026 ACM Conference on
Fairness, Accountability, and Transparency} (FAccT '26), pp.\ 1--44.
\url{https://doi.org/10.1145/3805689.3806763}

\bibitem[OpenAI(2025)]{openai2025modelspec}
OpenAI. \emph{The OpenAI Model Spec} (2025-09-12 revision). September 2025.
\url{https://model-spec.openai.com/2025-09-12.html}

\bibitem[Valbuena(2026)]{x1xhlol2026repo}
L.~Valbuena (x1xhlol). \emph{system-prompts-and-models-of-ai-tools}
[GitHub repository, aggregating reconstructed system prompts for 25+
coding-agent products; 141k+ stars, 501 commits at time of access].
Accessed July 2026.
\url{https://github.com/x1xhlol/system-prompts-and-models-of-ai-tools}

\bibitem[xAI(2026)]{xaigrokrepo2026}
xAI. \emph{grok-prompts} [GitHub repository, official raw system prompts
for Grok across chat, X-bot, and API deployment surfaces]. Accessed July 2026.
\url{https://github.com/xai-org/grok-prompts}

\bibitem[xAI(2025)]{xaiincident2025}
xAI. Statement on unauthorized modification to the Grok response bot on X.
Posted to X, May 15, 2025.
\url{https://x.com/xai/status/1923183620606619649}

\bibitem[Karmaker Santu and Feng(2023)]{teler2023}
S.~K. Karmaker Santu and D.~Feng. TELeR: A General Taxonomy of LLM Prompts
for Benchmarking Complex Tasks. \emph{arXiv:2305.11430}. Accepted to
Findings of EMNLP 2023.

\bibitem[Sahoo et al.(2024)]{liu2024survey}
P.~Sahoo, A.~K. Singh, S.~Saha, V.~Jain, S.~Mondal, and A.~Chadha.
A Systematic Survey of Prompt Engineering in Large Language Models:
Techniques and Applications. \emph{arXiv:2402.07927}, 2024.

\bibitem[Tian et al.(2025)]{promptdefects2025}
H.~Tian, C.~Wang, B.~Yang, L.~Zhang, and Y.~Liu. A Taxonomy of Prompt
Defects in LLM Systems. \emph{arXiv:2509.14404}, 2025.

\bibitem[Perez and Ribeiro(2022)]{perez2022}
F.~Perez and I.~Ribeiro. Ignore Previous Prompt: Attack Techniques For
Language Models. \emph{arXiv:2211.09527}, 2022. Presented at the NeurIPS
2022 ML Safety Workshop.

\bibitem[Yang et al.(2026)]{promptleak2026}
Y.~Yang, C.~Fu, T.~Zhang, R.~Zeng, Q.~Li, T.~Du, Z.~Wang, S.~Ji, and
W.~Chen. Understanding and Mitigating Prompt Leaking Attacks in
Real-World LLM-Based Applications. \emph{arXiv:2606.18673}, 2026.

\bibitem[Zhuang et al.(2025)]{proxyprompt2025}
Z.~Zhuang, M.-I. Nicolae, H.-P. Wang, and M.~Fritz. ProxyPrompt: Securing
System Prompts against Prompt Extraction Attacks. \emph{arXiv:2505.11459}, 2025.

\bibitem[Sahu et al.(2026)]{encodingdefense2026}
A.~Sahu, D.~Samanta, and R.~Soosahabi. Automated Framework to Evaluate
and Harden LLM System Instructions against Encoding Attacks.
\emph{arXiv:2604.01039}, 2026.

\end{thebibliography}

\end{document}
